\subsubsection{统一证书轴、holonomy 账本与三纤维联合末端正规形}\label{subsec:bridge-rh-unified-terminal-conclusions}
在既有“二变量自对偶行列式—完成化—unitary slice 审计—可寻址零点桥接合同（\(\mathbf{Bridge^{RH}}\)）”框架（\S\ref{subsec:bridge-rh-axiom-upgrade}，假设 \ref{assump:bridge-rh}）之上，本小节在四类既有接口之上给出三条统一性终端结论：地址先于取值的空值语义（公理 \ref{ax:pommin-address-null} 与 \S\ref{sec:recursive-addressing}；并见 \S\ref{subsec:xi-pw-type-safety-null} 的 unitary slice 类型安全口径），递归地址化的 \(\sigma\)-代数闭包（命题 \ref{prop:recursive-addressing-refinement}），prime-register 残差提升与可逆化（定义 \ref{def:pom-fold-prime-lift}--命题 \ref{prop:pom-fold-prime-lift-injective}），以及单位圆面积相位紧化与绕数证书的公理化接口（附录 \ref{app:unit-circle-phase-gate}，尤其是定义 \ref{def:dirichlet-trifiber-injection}--命题 \ref{prop:dirichlet-trifiber-injection-audit}）。
目标并非复述既有结构，而是把 \S\ref{subsec:bridge-rh-axiom-upgrade} 中“隐藏盘内数据需外置为可审计轴”的缺口（注 \ref{rem:bridge-rh-open-gaps}）提升为可组合、可证书化且可回放的终端接口，并给出面向三纤维（值/相位/证书轴）的统一编译判据。
在此口径下，所谓“第四寄存器”障碍可被严格化为：当 holonomy 账本在闭包内不平凡时，三纤维的“末端一次投影”联合正规形无法成立，扩展保存（公理 \ref{axiom:unit-circle-phase-extension-preserving}）强迫引入额外记录负载；本节将该负载压缩为 prime-register 轴上的可审计不变量。
\par
以下三条结论均以“可直接纳入审计/编译接口的最终陈述”给出。

\paragraph{预备结构与记号（仅列接口）}
为便于直接引用，仅固定下述最小接口：
\begin{enumerate}
  \item \textbf{地址与空值（\(\mathsf{NULL}\)）}：地址 \(a\in\{0,1\}^k\) 给出柱事件 \(C_a\) 与二元投影 \(\ind_{C_a}(x)\in\{0,1\}\)。在地址未给定之前，读出不定义（语义空值 \(\mathsf{NULL}\)），且“不定义”不等同于取值 \(0\)（公理 \ref{ax:pommin-address-null}）。
  \item \textbf{递归地址化的 \(\sigma\)-代数闭包}：派生柱事件属于旧层事件的平移闭包所生成的 \(\sigma\)-代数，递归地址化仅执行柱集细化而不引入外加可测集（命题 \ref{prop:recursive-addressing-refinement}）。
  \item \textbf{扩展保存（信息守恒的严格化）}：对一般非单射观测 \(\pi:\Omega\to X\)，若要求总体信息不丢失，则存在记录轴 \(r:\Omega\to R\) 使扩展观测 \(\widetilde\pi:=(\pi,r)\) 为单射（公理 \ref{axiom:unit-circle-phase-extension-preserving}）。
  \item \textbf{prime-register 残差提升}：prime-register 纤维 \(\mathcal{P}=\NN^{(\mathbb{P})}\) 由有限支撑素数指数向量给出，值映射 \(N(\mathbf r)=\prod_{p}p^{\mathbf r(p)}\) 由算术基本定理给出同构（定义 \ref{def:pom-prime-fiber}）。对折叠正规化轨迹的原子事件序列进行素数轴编码得到残差 \(r_{\mathsf S}(\omega)\)，并定义提升 \(\widetilde{\Fold}_m^{\mathsf S}(\omega)=(\Fold_m(\omega),r_{\mathsf S}(\omega))\)；对任意确定性策略 \(\mathsf S\) 该提升为单射（定义 \ref{def:pom-fold-prime-lift}，命题 \ref{prop:pom-fold-prime-lift-injective}）。
  \item \textbf{单位圆面积相位紧化与绕数证书}：无参紧化门 \(\upsilon(x)=\pi^{-1}\arctan x\in(-\tfrac12,\tfrac12)\) 与 Cayley 嵌入 \(\iota(x)=\frac{1+\mathrm{i}x}{1-\mathrm{i}x}\in\TT\) 将无界实轴压缩到单位圆端点（附录 \ref{app:unit-circle-phase-gate}）。并沿结构搬运在相位盒 \(\mathbb{U}=(-\tfrac12,\tfrac12)\) 上定义与 \((\RR,+,\cdot)\) 严格同构的运算 \((\boxplus,\boxtimes)\)，从而把数值发散/溢出统一改写为相位趋近边界的端点问题（定义 \ref{def:unit-circle-phase-transport-arith}，定理 \ref{thm:unit-circle-phase-field-iso}，注 \ref{rem:unit-circle-phase-bounded}）。
  在 Dirichlet--Cayley 扩展中，为恢复角变量的周期折叠信息，绕数证书 \(\kappa_\lambda(s)\in\ZZ\) 给出一致性约束 \(t\log\lambda+\mathrm{Arg}(z)=2\pi\kappa_\lambda(s)\in 2\pi\ZZ\)，并由三纤维注入 \(\mathcal{E}_\lambda(\sigma,t)=(|z|,u_{\mathrm C}(t),\kappa_\lambda(\sigma+\mathrm{i}t))\) 组织为可审计接口（定义 \ref{def:dirichlet-trifiber-injection}，命题 \ref{prop:dirichlet-trifiber-injection-audit}）。
\end{enumerate}

\begin{theorem}[统一性终端结论 A：prime-register 的统一证书轴（残差融合与单轴可逆化）]\label{thm:prime-register-unified-certificate-axis}
设 \(\pi:\Omega\to X\) 为任意观测，\(r:\Omega\to\mathcal{P}\) 为 prime-register 证书轴，\(\kappa:\Omega\to\ZZ\) 为整数证书轴。
若扩展观测
\[
(\pi,r,\kappa):\Omega\longrightarrow X\times\mathcal{P}\times\ZZ
\]
为单射，则存在一个\emph{固定}注入编码
\[
\mathrm{Enc}:\mathcal{P}\times\ZZ\hookrightarrow \mathcal{P}
\]
使得单一 prime-register 轴上的扩展观测
\[
\Omega\longrightarrow X\times\mathcal{P},\qquad
\omega\longmapsto\bigl(\pi(\omega),\mathrm{Enc}(r(\omega),\kappa(\omega))\bigr)
\]
仍为单射。
从而在“扩展保存”的严格口径下，离散折叠残差（prime-register）与整数型证书（例如绕数证书）可在同一 prime-register 轴内实现统一承载。
\end{theorem}
\begin{proof}
取素数序列 \(p_1<p_2<\cdots\)（其中 \(p_1=2\)）。
定义 prime-register 的轴重标定嵌入 \(\mathrm{sh}:\mathcal{P}\hookrightarrow\mathcal{P}\) 为
\[
(\mathrm{sh}(\mathbf r))(p_{t+1})=\mathbf r(p_t)\ (t\ge 1),
\qquad
(\mathrm{sh}(\mathbf r))(p_1)=0,
\]
并令其余坐标为 \(0\)。
同时取一固定注入 \(b:\ZZ\hookrightarrow\NN\)，例如
\[
b(k)=
\begin{cases}
2k, & k\ge 0,\\
-2k-1, & k<0.
\end{cases}
\]
定义
\(
\mathrm{Enc}(\mathbf r,k)\in\mathcal{P}
\)
为：在 \(p_1\)-坐标置 \(\mathrm{Enc}(\mathbf r,k)(p_1)=b(k)\)，在其余坐标置 \(\mathrm{Enc}(\mathbf r,k)(p)=\mathrm{sh}(\mathbf r)(p)\)。
则由 \(p_1\)-坐标可唯一恢复 \(k\)，并由其余坐标（将 \(p_{t+1}\) 反移回 \(p_t\)）唯一恢复 \(\mathbf r\)，故 \(\mathrm{Enc}\) 为注入。
\par
若 \(\bigl(\pi(\omega_1),\mathrm{Enc}(r(\omega_1),\kappa(\omega_1))\bigr)=\bigl(\pi(\omega_2),\mathrm{Enc}(r(\omega_2),\kappa(\omega_2))\bigr)\)，
则 \(\pi(\omega_1)=\pi(\omega_2)\) 且由 \(\mathrm{Enc}\) 的注入性得 \(r(\omega_1)=r(\omega_2)\)、\(\kappa(\omega_1)=\kappa(\omega_2)\)；
再由 \((\pi,r,\kappa)\) 的单射性推出 \(\omega_1=\omega_2\)。
证毕。
\end{proof}

\begin{remark}[意义（接口化表述）]\label{rem:prime-register-as-unified-certificate-axis}
定理 \ref{thm:prime-register-unified-certificate-axis} 把“扩展保存必须外置记录轴”（公理 \ref{axiom:unit-circle-phase-extension-preserving}）具体化为：\emph{prime-register 纤维 \(\mathcal{P}\) 可作为统一、可组合、可回放的通用证书轴}。
在该口径下，不同类型的不可逆性（折叠的多对一、周期角变量的同余折叠等）不再需要分属互异的记录域，而可在同一 prime-register 轴内实现单射化与审计闭环。
\end{remark}

\begin{theorem}[统一性终端结论 B：Holonomy--寄存器同型（证书路径与绕数证书的统一账本）]\label{thm:holonomy-register-isomorphism}
在投影词 \(2\)-范畴 \(\mathbf{PW}\) 的可审计片段中（定义 \ref{def:pom-projword-2cat}；并见注 \ref{rem:pom-projword-lift-proj-normalizer-demo} 与注 \ref{rem:pom-holonomy-cocycle-audit}），设每一条关键交换均被记录为可组合的交换证书，从而任何规约证书路径都可写为有限原子证书步序列 \(\gamma=(e_1,\dots,e_T)\)。
固定一组可审计注入编码 \(\mathrm{code}:\{e\}\hookrightarrow\NN_{\ge 1}\)（与定义 \ref{def:pom-fold-prime-lift} 同型），并令 \(p_1<p_2<\cdots\) 为素数序列。
定义 holonomy 账本映射
\[
\mathcal{H}:\{\text{证书路径}\}\longrightarrow \mathcal{P},
\qquad
\mathcal{H}(\gamma)(p_t):=\mathrm{code}(e_t)\ (1\le t\le T),
\]
其余坐标取 \(0\)。
则：
\begin{enumerate}
  \item \textbf{（可组合律）} 令 \(\mathrm{Sh}_k:\mathcal{P}\to\mathcal{P}\) 为素数轴移位算子 \((\mathrm{Sh}_k\mathbf r)(p_{t+k})=\mathbf r(p_t)\)（其余为 \(0\)），并对 \(\mathcal{H}\) 的值域定义拼接乘法
  \[
  \mathbf r_2\odot \mathbf r_1:=\mathbf r_1+\mathrm{Sh}_{T(\mathbf r_1)}(\mathbf r_2),
  \qquad
  T(\mathbf r):=\max\{t:\mathbf r(p_t)\neq 0\}\ (\text{约定 }T(0)=0).
  \]
  则对可纵向复合的证书路径 \(\gamma_1,\gamma_2\) 有
  \[
  \mathcal{H}(\gamma_2\circ\gamma_1)=\mathcal{H}(\gamma_2)\odot \mathcal{H}(\gamma_1).
  \]
  \item \textbf{（折叠残差为特例）} 当 \(\gamma\) 取为折叠正规化轨迹的原子重写事件序列时，\(\mathcal{H}(\gamma)\) 与 prime-register 残差 \(r_{\mathsf S}(\omega)\) 同型一致（定义 \ref{def:pom-fold-prime-lift}），并因此承继其“回放可逆/单射化”的性质（命题 \ref{prop:pom-fold-prime-lift-injective}）。
  \item \textbf{（绕数证书为特例）} 当 \(\gamma\) 对应于相位角变量的周期折叠一致性（定义 \ref{def:dirichlet-trifiber-injection}），则绕数证书 \(\kappa_\lambda(s)\in\ZZ\) 可由定理 \ref{thm:prime-register-unified-certificate-axis} 的 \(\mathrm{Enc}\) 编入 prime-register 轴，且一致性约束 \(t\log\lambda+\mathrm{Arg}(z)=2\pi\kappa_\lambda(s)\)（式 \eqref{eq:dirichlet-arg-consistency}）成为 \(\mathcal{P}\)-账本上的可审计整数条件。
  \item \textbf{（可观测等价的完备判据：在可审计片段内）} 若规约器固定且每一步关键交换均外置为可组合证书步（从而 \(\mathcal{H}\) 可回放），则任意两条实现链在外表读出一致时，其“路径依赖/不可交换残差”完全由 \(\mathcal{H}\) 的值刻画：同一外表读出且 \(\mathcal{H}\) 一致 \(\Longleftrightarrow\) 同一接口正规形与同一证书账本，从而在该片段上恢复为可判等价的自然同构类（参见定理 \ref{thm:pom-sem-conservative-on-slice} 的保守性口径）。
\end{enumerate}
\end{theorem}

\begin{theorem}[统一性终端结论 C：三纤维平坦性与“单次末端投影”联合正规形（地址--守恒--紧化的统一编译判据）]\label{thm:trifiber-flatness-terminal-projection-normal-form}
考虑同时涉及三类纤维读出门的实现链：值纤维上的折叠/条件期望投影 \(E_m\)（定义 \ref{def:pom-value-delay-rewrite}），相位纤维上的单位圆规范化读出（推论 \ref{cor:phase-projection-gate-budget}），以及为满足扩展保存而引入的 prime-register 证书轴（命题 \ref{prop:pom-fold-prime-lift-injective} 与定理 \ref{thm:prime-register-unified-certificate-axis}）。
若存在一条可核对的\textbf{联合纤维不可见/平坦性证书}，使得中间演化段对三类纤维（值纤维、相位纤维、证书纤维）均不可区分，则整条实现链可在 \(\mathbf{PW}\) 的可审计片段内等价规约为\textbf{单次末端投影的联合正规形}：所有中间出现的值投影与相位投影均可被推迟并融合为末端至多一次调用，而不会改变外表读出及其可核验性。
\par
\textbf{逆否判据（不可规约性与 \(\mathsf{NULL}\)）}：若联合 holonomy 账本（定理 \ref{thm:holonomy-register-isomorphism} 的 \(\mathcal{H}\)）在闭包内为非平凡类，则上述联合平坦性证书不存在；在此情况下，任何试图强制“无新增寄存器且仅保留纯相位可见域”的实现都会破坏扩展观测的单射化要求（公理 \ref{axiom:unit-circle-phase-extension-preserving}），从而在“地址先于取值”的语义中必然表现为不可寻址（读出不定义，即 \(\mathsf{NULL}\)，公理 \ref{ax:pommin-address-null}）。
\end{theorem}

\begin{remark}[闭环总结（终端接口）]\label{rem:bridge-rh-terminal-loop-summary}
以上三条终端结论共同给出一个可审计闭环：
\[
(\text{地址的 \(\sigma\)-代数闭包})
\ \perp\
(\text{扩展保存的单射化要求})
\ \perp\
(\text{单位圆相位紧化的周期一致性})
\]
三者在“prime-register 统一证书轴 + holonomy 账本 + 三纤维联合末端正规形”之下实现同一化。
在 \(\mathbf{Bridge^{RH}}\) 口径中，这一闭环把“隐藏盘内数据是否可被强制外置并在闭包内核验”（注 \ref{rem:bridge-rh-open-gaps}）转化为：是否存在可核对的 \(\mathcal{H}\)-平凡证书与相应的末端一次投影正规形；若不存在，则其失败必以 \(\mathsf{NULL}\) 与寄存器预算超界的形式显式外露。
\end{remark}

